\noindent Métodos baseados em Kernel são técnicas extremamente flexíveis que podem ser utilizados
para estender modelos ou algoritmos para regiões de decisão não lineares \cite{mohri2018foundations}. Estes métodos
são baseados nas conheecidas funções de kernel, os quais, sob condições de simetria e positiva definida, definem
produtos internos em espaços de alta dimensão, conhecidos como espaços de Hilbert. Dessa forma, este capítulo
terá como objetivo definir e apresentar os conceitos fundamentais sobre métodos baseados em kernel,
incluindo a definição de funções de kernel, o teorema do representer e Reproducing Kernel Hilbert
Space (RKHS).

\section{Kernels}

Por definição, uma função $K: \mathcal{X} \times \mathcal{X} \to \mathbb{R}$ é chamada de kernel sobre
$\mathcal{X}$.
A ideia central é construir um kernel $K$ tal que, para quaisquer dois pontos $x, x' \in \mathcal{X}$, $K(x, x')$
corresponda ao produto interno entre duas representações desses pontos em um espaço de Hilbert
$\mathcal{H}$, isto é,
$$
\forall x, x' \in \mathcal{X}, \quad K(x,x') = \langle \phi(x), \phi(x') \rangle,
$$
para algum mapeamento $\phi: \mathcal{X} \to \mathcal{H}$. O espaço $\mathcal{H}$ é chamado de
feature space.

Sob essa perspectiva, o kernel pode ser interpretado como uma medida de similaridade entre elementos
de $\mathcal{X}$, uma vez que o produto interno mensura a similaridade entre dois vetores.

Uma vantagem fundamental dessa abordagem é computacional. Isto é, muitas vezes é significativamente mais eficiente
calcular diretamente $K(x,x')$ do que computar explicitamente $\phi(x)$ e $\phi(x')$ e então avaliar
seu produto interno.

\subsection{Kernels positivos definidos}

Para que uma função $K$ possa ser utilizada como kernel, é necessário que sejam satisfeitas duas condições, que
a função seja simétrica e positiva definida. Satisfazendo essas condições, o kernel é capaz de induzir uma geometria
em um espaço de Hilbert.

Formalmente, um kernel $K: \mathcal{X} \times \mathcal{X} \to \mathbb{R}$ é dito simétrico e positivo
definido se, para qualquer conjunto finito $\{x_1, \ldots, x_n\} \subset \mathcal{X}$, a matriz de Gram
$$
\mathbf{K} = [K(x_i,x_j)]_{i,j=1}^n \in \mathbb{R}^{n \times n}
$$
for positiva semidefinida. Isto é, para qualquer vetor $\mathbf{c}=(c_1,\ldots,c_n)^\top \in \mathbb{R}^n$,
$$
\mathbf{c}^\top \mathbf{K}\mathbf{c}
=
\sum_{i=1}^n \sum_{j=1}^n c_i c_j K(x_i,x_j)
\geq 0.
$$

Equivalentemente, todos os autovalores de $\mathbf{K}$ são não negativos. Essa propriedade é essencial,
pois garante a existência de uma representação geométrica consistente do kernel.

\subsection{Espaços de Hilbert}

A noção de espaço de Hilbert fornece a estrutura matemática necessária para a teoria de kernels.

Por definição, um espaço de Hilbert $\mathcal{H}$ é um espaço vetorial munido de um produto interno
$\langle \cdot, \cdot \rangle_{\mathcal{H}}$ e que é completo com relação à norma induzida por esse produto interno. A norma em $\mathcal{H}$ é definida por
$$
\forall h \in \mathcal{H}, \quad \|h\|_{\mathcal{H}} = \sqrt{\langle h,h\rangle}.
$$

A completude implica que toda sequência de Cauchy em $\mathcal{H}$ converge para um elemento do próprio
espaço, garantindo estabilidade analítica para problemas de otimização e aproximação.

No contexto de kernels, o espaço de Hilbert atua como o espaço onde os dados são representados via o
mapeamento $\phi$.

\subsection{Reproducing Kernel Hilbert Spaces}

A relação entre kernels positivos definidos e espaços de Hilbert é formalizada pelo teorema.

Seja $K: \mathcal{X} \times \mathcal{X} \to \mathbb{R}$ um kernel simétrico e positivo definido. Então existe um espaço de Hilbert $\mathcal{H}$ e um mapeamento $\phi: \mathcal{X} \to \mathcal{H}$ tais que

$$
\forall x,x' \in \mathcal{X}, \quad K(x,x') = \langle \phi(x), \phi(x') \rangle_{\mathcal{H}}.
$$

Além disso, esse espaço satisfaz a chamada propriedade de reprodução:

$$
\forall h \in \mathcal{H}, \forall x \in \mathcal{X}, \quad h(x) = \langle h, K(x,\cdot)\rangle_{\mathcal{H}}.
$$

Essa propriedade estabelece que a avaliação funcional é contínua e pode ser representada por meio do próprio
 kernel.

O espaço $\mathcal{H}$ associado a $K$ é chamado de Reproducing Kernel Hilbert Space (RKHS), desempenhando
papel central em métodos de aprendizado estatístico, particularmente em problemas de regularização,
estimação não paramétrica e aprendizado supervisionado.

\subsection{Teorema do Representante}

O resultado mais interessante para a definição do modelo GG-KM é o Teorema do Representante. Sua principal
importância reside no fato de garantir que, embora a otimização ocorra em um espaço potencialmente infinito
-dimensional, a solução ótima pode ser expressa em termos de uma combinação finita de kernels avaliados nas
observações da amostra.

Formalmente, seja $K \text{ : } \mathcal{X} \times \mathcal{X} \to \mathbb{R}$ um kernel positivo definido e
$\mathcal{H}$ seu correspondente RKHS. Assim, para qualquer função $G \text{ : } \mathbb{R} \to \mathbb{R}$ não
decrescente e qualquer função perda $L \text{ : } \mathbb{R}^n \to \mathbb{R} \cup \{+\infty\}$, o problema de otimização
$$
{\arg\min}_{h\in\mathcal{H}} F\left(h\right) = {\arg\min}_{h\in\mathcal{H}} G\left(\|h\|_{\mathcal{H}}\right) + L\left(h\left(x_1\right), \ldots, h\left(x_n\right)\right)
$$
admite uma solução na forma $h^{*} = \sum_{i=1}^{n}\alpha_i K\left(x_i, \cdot\right)$. Se $G$ é assumido como sendo crescente,
então qualquer solução tem esta forma.


Esse resultado é particularmente relevante porque reduz um problema de otimização funcional em dimensão
infinita para um problema de otimização em dimensão finita, dependente apenas dos coeficientes
$\boldsymbol{\alpha} = (\alpha_1, \ldots, \alpha_n)$. Na prática, isso torna viável o uso de métodos
baseados em kernel, além de evidenciar que toda a informação necessária para a estimação está contida na
matriz kernel $K$, sem necessidade de explicitar o mapeamento $\phi(\cdot)$.

% \subsection{Principais funções kernel}

% Uma vez estabelecido pelo Teorema do Representante que a solução ótima pode ser escrita
% como uma combinação linear de kernels avaliados nos dados observados, a escolha da função
% kernel passa a desempenhar papel central na modelagem. Diferentes kernels induzem distintas
% geometrias no espaço de características, o que afeta diretamente a flexibilidade e a capacidade
% de generalização do modelo.

% Entre os kernels mais utilizados, destacam-se:

% \begin{itemize}

%     \item \textbf{Kernel linear}
%     \[
%     K(x_i,x_j)=x_i^\top x_j
%     \]

%     \item \textbf{Kernel polinomial}
%     \[
%     K(x_i,x_j)=\left(x_i^\top x_j + c\right)^d
%     \]

%     \item \textbf{Kernel Gaussiano (RBF)}
%     \[
%     K(x_i,x_j)=\exp\left(-\gamma \|x_i-x_j\|^2\right)
%     \]

%     \item \textbf{Kernel Laplaciano}
%     \[
%     K(x_i,x_j)=\exp\left(-\gamma \|x_i-x_j\|\right)
%     \]

%     \item \textbf{Kernel Sigmoidal}
%     \[
%     K(x_i,x_j)=\tanh\left(\gamma x_i^\top x_j + c\right)
%     \]

%     \item \textbf{Kernel Cauchy}
%     \[
%     K(x_i,x_j)=\left(1+\frac{\|x_i-x_j\|^2}{\sigma^2}\right)^{-1}
%     \]

%     \item \textbf{Kernel exponencial}
%     \[
%     K(x_i,x_j)=\exp\left(-\gamma \|x_i-x_j\|\right)
%     \]
% \end{itemize}

A escolha do kernel determina implicitamente a classe de funções admissíveis no espaço RKHS e,
consequentemente, o grau de complexidade que o modelo pode representar. Neste contexto, os kernels podem ser
visto como um hiperparâmetro de complexidade do modelo, bem como as variáveis que os compõe, $\gamma\text{, } d\text{ e } \sigma$.

\section{GG-KM}

O modelo proposto GG-KM, nada mais é que uma proposta de modelar o valor esperado dos riscos latentes do
modelo de tempo de promoção a partir da integração de kernels, de forma que seja uma extensão semi-paramétrica.
Esta seção terá como objetivo descrever o processo de estimação e fundamentação do modelo.

\subsection{Extensão não paramétrica via kernel}

Embora o modelo de tempo de promoção clássico permita incorporar covariáveis por meio de funções paramétricas
para $\theta(\mathbf{x})$, essa abordagem pode ser restritiva em cenários onde a relação entre as covariáveis
e a fração de cura apresenta estruturas complexas ou altamente não lineares.

Neste contexto, propõe-se modelar a intensidade média dos riscos latentes $\theta\left(\mathbf{x}\right)$ por
meio de uma função não paramétrica $f\left(\mathbf{x}\right) = \log\left(\theta\left(\mathbf{x}\right)\right)>0$.
Dessa forma, suponha que a estimação do modelo de promoção de cura seja feita por meio da maximização da log-verossimilhança
penalizada:
\begin{equation}
\label{eq:ggkm_penalized_likelihood}
\max_{f} \sum_{i=1}^n l(\mathbf{x}_i) - \lambda \|f\|^2 = \min_{f} \sum_{i=1}^n - l(\mathbf{x}_i) + \lambda \|f\|^2
\end{equation}
em que $l(\mathbf{x}_i)$ representa a log-verossimilhança e $\lambda$ é uma parâmetro de regularização da complexidade
do modelo. No entanto, supondo que a função $f$ pertença a um Reproducing Kernel Hilbert Space (RKHS)
$\mathcal{H}_K$, associado a um kernel positivo definido $K(\cdot,\cdot)$, o Teorema do Representante
garante que a solução do problema da Equação \ref{eq:ggkm_penalized_likelihood} pode ser expressa da
seguinte forma:
\begin{equation}
f^*(\mathbf{x}) = \sum_{i=1}^{n}\alpha_i K(\mathbf{x},\mathbf{x}_i),
\end{equation}
em que $\alpha_1,\ldots,\alpha_n \in \mathbb{R}$ são coeficientes a serem estimados. Consequentemente, a
intensidade média dos riscos latentes para o indivíduo $i$ pode ser escrita como
\begin{equation}
\theta(\mathbf{x}_i) = \exp\left(\sum_{j=1}^{n}\alpha_j K(\mathbf{x}_i,\mathbf{x}_j)\right).
\end{equation}

% Para o modelo de tempo de promoção, tem-se que a função log-verossimilhança é dada por:
% \begin{equation}
% \label{eq:ggkm_log_likelihood}
% l(\mathbf{x}_i) = \delta_i \log\left(h_{\text{pop}}\left(t_i|\mathbf{x}_i\right)\right) + \log\left(S_{\text{pop}}\left(t_i|\mathbf{x}_i\right)\right)\text{.}
% \end{equation}

\subsection{Estimação via algoritmo EM penalizado}

A representação obtida pelo Teorema do Representante permite reescrever o modelo de tempo de promoção em
termos de um número finito de parâmetros. Seja
\begin{equation}
f(\mathbf{x}_i) = \sum_{j=1}^{n} \alpha_j K(\mathbf{x}_i,\mathbf{x}_j),
\end{equation}

Sob a formulação do modelo de tempo de promoção, assume-se que o número de riscos latentes associados ao
indivíduo $i$ segue uma distribuição de Poisson com média $\theta(\mathbf{x}_i)$, isto é,
\begin{equation}
N_i \mid \mathbf{x}_i
\sim
\text{Poisson}
\left(
\theta(\mathbf{x}_i)
\right).
\end{equation}

Entretanto, os riscos competitivos $N_i$ associados ao evento de interesse não são observados. Assim, a
estimação dos parâmetros pode ser conduzida por meio do algoritmo EM. Neste caso, considera-se o algoritmo
EM penalizado \cite{mclachlan2008algorithm}.

Dessa forma, considere os dados completos $D_{\text{comp}} = \left\{t_i, \delta_i, N_i, \mathbf{x}_i\right\}^n_{i=1}$,
em que $t_i$ representa o tempo observado, $\delta_i$ o indicador de falha e $N_i$ o número de riscos latentes
associados ao $i\text{-ésimo}$ indíviduo. A log-verossimilhança completa, menos uma constante aditiva, pode ser escrita
como
\begin{equation}
\begin{aligned}
l_c
=
&
\sum_{i=1}^{n}
\left[
-\exp\!\bigl(f(\mathbf{x}_i)\bigr)
+
N_i f(\mathbf{x}_i)
\right]
\\
&
+
\sum_{i=1}^{n}
\left[
\delta_i
\log f_{GG}(t_i;a,d,p)
\right]
\\
&
+
\sum_{i=1}^{n}
\left[
(N_i-\delta_i)
\log S_{GG}(t_i;a,d,p)
\right].
\end{aligned}
\end{equation}

Para controlar a complexidade da função estimada, adiciona-se uma penalização baseada na norma do RKHS. Dessa forma, obtém-se a log-verossimilhança completa penalizada
\begin{equation}
\ell_{c, P} = l_c - \frac{\lambda}{2} \boldsymbol{\alpha}^{\top} K \boldsymbol{\alpha},
\end{equation}
em que
\begin{equation}
\boldsymbol{\alpha} = (\alpha_1,\ldots,\alpha_n)^{\top}
\end{equation}
e $K$ denota a matriz de kernel com elementos $K_{ij}=K(\mathbf{x}_i,\mathbf{x}_j)$.

O algoritmo EM consiste em alternar entre o cálculo da esperança condicional dos riscos latentes e a
maximização da esperança da log-verossimilhança completa penalizada.

No passo E, calcula-se
\begin{equation}
Q\!\left(
\Theta \mid \Theta^{(k)}
\right)
=
\mathbb{E}
\left[
\ell_{c,P}
\mid
\mathcal{D}_{\mathrm{obs}},
\Theta^{(k)}
\right].
\end{equation}
em que $\Theta = \left(\boldsymbol{\alpha},a,d,p\right)$.

Utilizando propriedades do modelo de tempo de promoção, obtém-se que

\begin{equation}
\widehat{N}_i^{(k+1)}
=
\mathbb{E}
\!\left[
N_i
\mid
\mathcal{D}_{obs},
\Theta^{(k)}
\right]
=
\delta_i
+
\theta_i^{(k)}
S_{GG}
\!\left(
t_i;
a^{(k)},
d^{(k)},
p^{(k)}
\right),
\end{equation}
em que
\begin{equation}
\theta_i^{(k)}
=
\exp
\!\left(
f^{(k)}(\mathbf{x}_i)
\right)
=
\exp
\!\left(
\sum_{j=1}^{n}
\alpha_j^{(k)}
K(\mathbf{x}_i,\mathbf{x}_j)
\right).
\end{equation}

Substituindo as quantidades não observadas por suas respectivas
esperanças condicionais, a função auxiliar penalizada assume a forma
\begin{equation}
Q_P
=
Q_{\alpha}
+
Q_{GG},
\end{equation}
com
\begin{equation}
Q_{\alpha}
=
\sum_{i=1}^{n}
\left[
-\exp\!\bigl(f(\mathbf{x}_i)\bigr)
+
\widehat N_i^{(k+1)}
f(\mathbf{x}_i)
\right]
-
\frac{\lambda}{2}
\boldsymbol{\alpha}^{\top}
K
\boldsymbol{\alpha},
\end{equation}
e
\begin{equation}
Q_{GG}
=
\sum_{i=1}^{n}
\left[
\delta_i
\log f_{GG}(t_i;a,d,p)
+
\bigl(\widehat N_i^{(k+1)}-\delta_i\bigr)
\log S_{GG}(t_i;a,d,p)
\right].
\end{equation}

Observa-se que a função auxiliar penalizada é separável em dois blocos de parâmetros,
\(\boldsymbol{\alpha}\) e \((a,d,p)\). Consequentemente, o passo M pode ser realizado por
otimização alternada.

Inicialmente, atualizam-se os coeficientes do kernel resolvendo
\begin{equation}
\boldsymbol{\alpha}^{(k+1)}
=
\arg\max_{\boldsymbol{\alpha}}
Q_{\alpha}.
\end{equation}

Em seguida, atualizam-se os parâmetros da distribuição Gama Generalizada por meio de
\begin{equation}
(a^{(k+1)},d^{(k+1)},p^{(k+1)})
=
\arg\max_{a,d,p}
Q_{GG}.
\end{equation}

A otimização de ambos problemas de otimização é realizada numericamente através de L-BFGS-B, utilizando-se
as derivadas analíticas das respectivas funções objetivos.

O processo iterativo é repetido até a convergência da log-verossimilhança penalizada observada,
\begin{equation}
\ell_P
=
\sum_{i=1}^{n}
\left[
\delta_i
\left(
f(\mathbf{x}_i)
+
\log f_{GG}(t_i;a,d,p)
\right)
-
\exp\!\bigl(f(\mathbf{x}_i)\bigr)
F_{GG}(t_i;a,d,p)
\right]
-
\frac{\lambda}{2}
\boldsymbol{\alpha}^{\top}
K
\boldsymbol{\alpha},
\end{equation}
sendo adotado como critério de parada
\begin{equation}
\left|
\ell_P^{(k+1)}
-
\ell_P^{(k)}
\right|
<
\varepsilon,
\end{equation}
em que \(\varepsilon > 0\) representa uma tolerância previamente especificada.

\begin{algorithm}[H]
\caption{Algoritmo EM penalizado para estimar o modelo GG-KM}
\begin{algorithmic}[1]

\State Inicialize
$
\alpha^{(0)},\,
a^{(0)},\,
d^{(0)},\,
p^{(0)}
$
e escolha uma tolerância $\varepsilon>0$.

\State Defina $k \gets 0$.

\Repeat

\State Calcule
\[
\theta_i^{(k)}
=
\exp\!\left(
\sum_{j=1}^{n}
\alpha_j^{(k)}
K(x_i,x_j)
\right),
\qquad i=1,\ldots,n.
\]

\State Calcule
\[
F_i^{(k)}
=
F_{GG}
\!\left(
t_i;
a^{(k)},
d^{(k)},
p^{(k)}
\right),
\qquad
S_i^{(k)}
=
1-F_i^{(k)}.
\]

\State \textbf{E-step:} para $i=1,\ldots,n$, calcule
\[
\widehat N_i^{(k+1)}
=
\delta_i
+
\theta_i^{(k)}
S_i^{(k)}.
\]

\State \textbf{M-step ($\boldsymbol{\alpha}$):}
\[
\alpha^{(k+1)}
=
\arg\max_{\alpha}
\left\{
\sum_{i=1}^{n}
\left[
-\theta_i
+
\widehat N_i^{(k+1)}
\log(\theta_i)
\right]
-
\frac{\lambda_{\mathrm{reg}}}{2}
\alpha^\top K\alpha
\right\},
\]

onde

\[
\theta_i
=
\exp
\left(
\sum_{j=1}^{n}
\alpha_j
K(x_i,x_j)
\right).
\]

\State \textbf{M-step (GG):}
\[
(a^{(k+1)},d^{(k+1)},p^{(k+1)})
=
\arg\max_{a,d,p}
\sum_{i=1}^{n}
\left[
\delta_i
\log f_{GG}(t_i;a,d,p)
+
(\widehat N_i^{(k+1)}-\delta_i)
\log S_{GG}(t_i;a,d,p)
\right].
\]

\State Calcule

\[
f_i^{(k+1)}
=
\sum_{j=1}^{n}
\alpha_j^{(k+1)}
K(x_i,x_j),
\qquad
\theta_i^{(k+1)}
=
\exp\!\left(f_i^{(k+1)}\right).
\]

\State Calcule a log-verossimilhança penalizada observada

\[
\begin{aligned}
\ell_P^{(k+1)}
=
&
\sum_{i=1}^{n}
\Bigg[
\delta_i
\Bigg(
f_i^{(k+1)}
+
\log
f_{GG}
\!\left(
t_i;
a^{(k+1)},
d^{(k+1)},
p^{(k+1)}
\right)
\Bigg)
\\
&
\qquad
-
\theta_i^{(k+1)}
F_{GG}
\!\left(
t_i;
a^{(k+1)},
d^{(k+1)},
p^{(k+1)}
\right)
\Bigg]
\\
&
-
\frac{\lambda_{\mathrm{reg}}}{2}
(\alpha^{(k+1)})^\top
K
\alpha^{(k+1)}.
\end{aligned}
\]

\State Calcule

\[
\Delta^{(k+1)}
=
\left|
\ell_P^{(k+1)}
-
\ell_P^{(k)}
\right|.
\]

\State $k \gets k+1$.

\Until{$\Delta^{(k+1)}<\varepsilon$}

\end{algorithmic}
\end{algorithm}
